% Converted from README.md
\subsection{Korišćenje C bibliotečkih funkcija i simboličkog okruženja}\label{koriux161ux107enje-c-biblioteux10dkih-funkcija-i-simboliux10dkog-okruux17eenja}

Kada se KLEE koristi nad programom koji poziva C bibliotečke funkcije, potrebno je navesti opcije \texttt{-{}-libc=uclibc -{}-posix-runtime}. One omogućavaju ispravno linkovanje i rad sa simboličkim okruženjem. Opcija \texttt{-{}-allow-external-sym-calls} dozvoljava pozivanje spoljnih funkcija sa simboličkim vrednostima.

Dodatne opcije za rad sa simboličkim okruženjem:

\begin{enumerate}
\def\labelenumi{\arabic{enumi}.}
\tightlist
\item
  \texttt{-{}-sym-arg\ \textless{}N\textgreater{}} - simbolički argument dužine najviše N ~
\item
  \texttt{-{}-sym-args\ \textless{}MIN\textgreater{}\ \textless{}MAX\textgreater{}\ \textless{}N\textgreater{}} - najmanje MIN, a najviše MAX simboličkih argumenata komandne linije, svi maksimalne dužine N ~
\item
  \texttt{-{}-sym-files\ \textless{}NUM\textgreater{}\ \textless{}N\textgreater{}} - NUM simboličkih datoteka veličine N (ne uključujući stdin) ~
\item
  \texttt{-{}-sym-stdin\ \textless{}N\textgreater{}} - simbolički standardni ulaz veličine N ~
\item
  \texttt{-{}-sym-stdout} - simbolički standardni izlaz
\end{enumerate}

Korisna je i opcija \texttt{-{}-optimize}, koja uklanja mrtav kod tokom istraživanja putanja i može da ubrza izvršavanje. Svaki poziv eksterne funkcije usložnjava pronalaženje putanja i generisanje test primera.

\subsubsection{Primer sa simboličkim argumentom komandne linije}\label{primer-sa-simboliux10dkim-argumentom-komandne-linije}

Primer \texttt{04\_libc} sadrži program za validaciju lozinke jednog sistema koja se zadaje kao argument komandne linije i za koju jedino znamo da je maksimalne dužine \texttt{100}. Pretpostaviti da postoji samo jedna valjana kombinacija za lozinku. Zadatak je pronaći lozinku bez rešavanja zadatih jednačina koje su postavljene u programu.

Zadatak ćemo rešiti označavanjem putanje koja nas zanima pozivom funkcije \texttt{klee\_assert(0)}, a to je upravo nakon ispisa poruke o uspešno pronađenoj lozinki (datoteka \texttt{password\_sym.c}). Nakon što se generiše bitkod \texttt{password\_sym.bc}, dodajemo potrebne opcije prilikom poziva KLEE alata:

\begin{minted}{bash}
$ klee --optimize --libc=uclibc --posix-runtime password_sym.bc -sym-arg 100
\end{minted}

Nakon što je prošlo par sekundi, može se primetiti da je generisana samo jedna poruka o grešci i čak \texttt{102} putanje. Da bismo dobili onu koja nas zanima, izvući ćemo naziv datoteke ekstenzije \texttt{.err} i pročitati odgovarajuću \texttt{.ktest} datoteku:

\begin{minted}{bash}
$ ls klee-last/*err
klee-last/test000026.assert.err
$ ktest-tool test000026.ktest
\end{minted}

\begin{minted}{text}
ktest file : 'test000026.ktest'
args       : ['password_sym.bc', '-sym-arg', '100']
num objects: 2
object    0: name: 'arg0'
object    0: size: 101
object    0: data: 'Pandi_panda\x00\xff\xff\xff\xff
\xff\xff\xff\xff\xff\xff\xff\xff\xff\xff\xff\xff\xff
\xff\xff\xff\xff\xff\xff\xff\xff\xff\xff\xff\xff\xff
\xff\xff\xff\xff\xff\xff\xff\xff\xff\xff\xff\xff\xff
\xff\xff\xff\xff\xff\xff\xff\xff\xff\xff\xff\xff\xff
\xff\xff\xff\xff\xff\xff\xff\xff\xff\xff\xff\xff\xff
\xff\xff\xff\xff\xff\xff\xff\xff\xff\xff\xff\xff\xff
\xff\xff\xff\xff\xff\xff\xff'
object    1: name: 'model_version'
object    1: size: 4
object    1: data: 1
\end{minted}

Dobili smo da je ulaz \texttt{Pandi\_panda} koji prati terminirajuća nula ispravna lozinka.
